% StateScope system-demonstration paper draft for AACL-IJCNLP 2026.
% Upload this file into the official ACL Overleaf project (which supplies acl.sty).
% AACL-IJCNLP 2026 demos are single-blind, so author names must be present.
\documentclass[11pt]{article}

\usepackage[final]{acl}
\usepackage{times}
\usepackage{latexsym}
\usepackage[T1]{fontenc}
\usepackage[utf8]{inputenc}
\usepackage{microtype}
\usepackage{inconsolata}
  \usepackage{graphicx}
  \usepackage{booktabs}
  \usepackage{multirow}
  \usepackage{xcolor}
  \usepackage{amsmath}
  \usepackage{amssymb}
  \usepackage{url}

% -----------------------------------------------------------------------------
% Replace every placeholder below before submission.  A public demo/package link
% and a <=2.5 minute video are mandatory for the AACL-IJCNLP 2026 demo track.
% -----------------------------------------------------------------------------
\newcommand{\systemname}{\textsc{StateScope}}
\newcommand{\codeurl}{https://github.com/USERNAME/DriftMath}
\newcommand{\demourl}{https://YOUR-LIVE-DEMO-OR-INSTALL-PAGE}
\newcommand{\videourl}{https://youtu.be/VIDEO-ID}
\newcommand{\softwarelicense}{[ADD SOFTWARE LICENSE, e.g., Apache-2.0]}

\title{\systemname: An Extensible Framework for Controlled,\\
Interactive Analysis of Mathematical Reasoning}

\author{First Author \\
  Department, Institution \\
  City, Country \\
  \texttt{first.author@example.edu}
  \And
  Second Author \\
  Department, Institution \\
  City, Country \\
  \texttt{second.author@example.edu}}

\begin{document}
\maketitle

\begin{abstract}
We present \systemname, an extensible framework for controlled, interactive analysis
of language-model mathematical reasoning.  The framework is organized around four
capabilities.  First, it generates seeded, fresh-parameter problem instances with
SymPy-constructed reference traces,
reducing dependence on static benchmarks and their contamination risk.  Second, it
controls reasoning context through typed operations and an explicit symbolic state,
rather than repeatedly exposing an unconstrained conversation history.  Third, it
supports multi-turn model execution and a deterministic counterfactual sandbox.
Fourth, core registries for problem families, model backends, and state systems,
together with a central operation-specification table, make the runtime
straightforward to extend.  As a flagship use case, \systemname{} detects
\emph{solution-state drift}: a variable binding, constraint, or intermediate value
diverges even when the final answer appears correct.  A browser interface compares
model-owned state with a CAS-verified external ledger and explains the first
divergence without an LLM judge.  Experiments on 250 dynamically generated problems
and four locally served open models validate the framework and illustrate the value
of explicit state ownership.  The system exports auditable traces for researchers,
developers, and educators.
Code, installable demo, and video are available at \url{\codeurl},
\url{\demourl}, and \url{\videourl}, respectively.
\end{abstract}

\section{Introduction}

Mathematical-reasoning research is often built around static datasets, uncontrolled
prompt histories, and a single final-answer metric.  This makes four practical
questions difficult to study.  Can evaluation instances be refreshed to reduce
memorization concerns?  Can the information available at each reasoning step be
specified exactly?  Can models and users interact with, inspect, and revise a live
reasoning process?  Can a new task family or inference system be added without
rewriting an evaluation stack?

\systemname{} is designed around these questions.  It combines a dynamic symbolic
problem generator, a context-controlled operation loop, an interactive replay UI,
and a registry-based research runtime.  Together, these components turn a benchmark
script into a reusable framework for constructing and studying stateful reasoning
systems.  The implementation underneath the interface is the DriftMath runtime;
\systemname{} is its end-to-end demonstration and analysis environment.

Our primary demonstration studies a limitation of final-answer evaluation.
Intermediate reasoning can be inconsistent even when cancellation, guessing, or a
later correction produces the right answer.  Process supervision and error-location
benchmarks make intermediate correctness explicit
\citep{lightman2023verify,zheng2025processbench}, but they typically score a fixed
natural-language solution.  They do not expose the evolving symbolic state that a
solver must carry across variable bindings, branch choices, domain constraints, and
reused lemmas.

We call divergence in that evolving state \emph{solution-state drift}.  For example,
after squaring $\sqrt{x+2}=x$, a solver must carry both the resulting candidates and
the original domain constraint; checking only the reported root cannot reveal when
one of these components was dropped.  The practical question is therefore not only
``Was the answer correct?'', but ``At which operation did the live state first cease
to match a verified derivation, what changed, and what happened downstream?''

The framework contributions are:
(1) \textbf{dynamic generation}: seeded, fresh-parameter, provenance-tracked
instances with SymPy-constructed gold traces;
(2) \textbf{context control}: a typed, one-operation-per-turn protocol whose next
prompt contains the problem and current symbolic state rather than an ever-growing
free-form transcript;
(3) \textbf{multi-turn interaction}: iterative model execution, paired side-by-side
system inspection, and a human-editable deterministic replay sandbox; and
(4) \textbf{extensibility}: core interfaces and registries for problem families,
model backends, and state-management systems, plus a single specification table for
the typed operation layer.
Solution-state drift detection is the unifying use case through which the paper
demonstrates these capabilities.  Unlike a learned critic, mathematical verdicts and
state differences and their explanations after generation are derived
deterministically; operation verdicts use SymPy where an operation is marked
CAS-checkable.

\begin{figure*}[t]
\centering
\setlength{\fboxsep}{5pt}
\small
\fbox{\parbox[c][1.05cm][c]{.15\textwidth}{\centering Problem\\+ local model}}
\hfill$\longrightarrow$\hfill
\fbox{\parbox[c][1.05cm][c]{.17\textwidth}{\centering Text--JSON adapter\\one typed operation}}
\hfill$\longrightarrow$\hfill
\fbox{\parbox[c][1.05cm][c]{.20\textwidth}{\centering System C: model claim\\System D: typed ledger}}
\hfill$\longrightarrow$\hfill
\fbox{\parbox[c][1.05cm][c]{.16\textwidth}{\centering SymPy verifier\\+ gold alignment}}
\hfill$\longrightarrow$\hfill
\fbox{\parbox[c][1.05cm][c]{.15\textwidth}{\centering Diff, replay,\\and export}}
\caption{\systemname{} framework.  Dynamic problems enter a context-controlled,
multi-turn operation loop.  Pluggable state systems feed a shared verifier and the
interactive inspection, replay, and export layer.}
\label{fig:architecture}
\end{figure*}

\section{Framework Design}

\subsection{Dynamic, contamination-resistant generation}

Instead of treating a static collection as the framework boundary, each problem
family implements a generator that samples fresh parameters from a seeded template
and constructs its symbolic solution trace with SymPy operations.  Unit tests check
the family and oracle behavior, but the current dataset builder does not run a
separate proof checker over every completed trace.  A manifest records the
seed, family counts, package version, source policy, and code revision.  Regeneration
therefore produces a reproducible evaluation set while allowing new surface
instances for subsequent studies.

This design is \emph{contamination-resistant}, not a claim of contamination proof.
A pretrained model may have encountered the underlying mathematical form, but it
cannot simply retrieve the freshly parameterized instance from a public test file.
New generators can implement fully synthetic tasks, licensed template
reinstantiation, or externally sourced tasks with an explicit contamination-risk
label; the headline demo uses only the synthetic path.

\subsection{Explicit context control}

The runtime exposes the model to a deliberately bounded view at every turn: the
problem statement, the current structured state, the allowed operation schemas, and
progress guardrails.  It does not resend a free-form transcript of prior rationales.
Consequently, an experimenter can identify what information is available, replace
the state-management policy, and measure failures in carrying state separately from
failures in choosing the next operation.  The ledger condition further makes
context state an executable system object rather than model-authored prose.

\subsection{Symbolic state and operations}

At reasoning step $t$, a state $S_t$ contains live and discharged bindings,
constraints, a current expression or equation, candidate solutions, dependency
nodes, and an optional final answer.  Expressions are serialized as strings but
compared semantically with SymPy \citep{meurer2017sympy}; equality is never based on
surface text.  Each problem also has a CAS-constructed gold trace
$G=(G_0,\ldots,G_{T-1})$.

The model emits exactly one JSON object per turn containing an operation name,
typed arguments, its complete claimed post-state, a completion flag, and a short
rationale.  The allowed operation set is family-specific and includes operations
such as \texttt{bind}, \texttt{square\_equation}, \texttt{filter\_candidates},
\texttt{establish\_lemma}, and \texttt{report}.  A bounded two-attempt repair loop
handles malformed JSON; validation rejects repeated bindings and premature
completion.  The same adapter and tool implementation are used in both conditions.

\subsection{State ownership comparison}

\paragraph{System C: tools + textual state.}
The operation is executed and checked on a scratch ledger, but the model's claimed
post-state is the trace state.  A stale or incomplete claim can therefore propagate
even when the selected operation itself is valid.

\paragraph{System D: tools + typed ledger.}
The model chooses the operation, while the runtime validates its arguments, applies
it to the ledger, and CAS-verifies the result.  The verified ledger snapshot becomes
the trace state.  Invalid or unverifiable operations stop the run rather than
silently changing state.  Thus, System D does not make the model a better planner;
it tests whether external state ownership prevents bookkeeping errors.

\subsection{Drift analysis}

For aligned step $t$, \systemname{} compares $S_t$ and $G_t$ component-wise.  The
corruption-onset depth (COD) is the first unequal step, and propagation length counts
later unequal aligned steps.  The UI identifies which of bindings, constraints,
current expression/equation, candidates, or final answer first diverged.  It also
shows whether the operation at COD passed CAS verification, separating a valid
operation with a stale carried state from an invalid transformation.  We use the
term \emph{drift} for divergence from the canonical oracle trajectory; a CAS-valid
alternative derivation is not automatically declared mathematically invalid.

The original benchmark metric, aligned state fidelity, is the fraction of aligned
steps whose post-states equal gold.  Because a short trace has fewer opportunities
to disagree, we additionally report strict gold-normalized fidelity:
\begin{equation}
 \mathrm{SF}_{g}=\frac{1}{T}\sum_{t=0}^{T-1}
 \mathbb{1}[t\ \text{is emitted and}\ S_t\equiv G_t].
 \label{eq:sfg}
\end{equation}
Missing steps therefore receive zero credit.  Final-answer accuracy is computed
independently by symbolic equality.

\subsection{Extensible research interfaces}

Problem families share a generator interface; systems share a solver interface;
model backends resolve through a model registry and YAML configuration; and all
operation descriptions, examples, validation rules, and tool schemas derive from a
single operation-specification table.  The browser obtains its model list from the
same AACL model catalog used by batch evaluation.  Adding another locally served
model that uses the existing backend requires a model configuration and one catalog
entry rather than core-logic changes.  Adding a mathematical
phenomenon requires a family generator, gold-trace builder, and any genuinely new
operation specifications.  Trace, metric, replay, and export components remain
unchanged.  The core runtime is extensible, although the current demo UI explicitly
hardcodes Systems C and D and the AACL model catalog.  Metrics and report analyses
are modular functions, not a plugin registry.  This separation is the central
framework-level selling point: the drift demo is one analysis built from reusable
execution primitives.

\section{Multi-Turn Interactive Demonstration}

The current UI exposes a curated set of generated problems.  The user selects one,
a model, and an execution route, then chooses either \emph{Run both} or a
one-turn-at-a-time \emph{Step through} view.  For a low-cost hosted presentation the
interface defaults to GPT-5 nano when an OpenAI credential is available; the
headline evaluation in Section~4 remains the preregistered open, locally served
model matrix.  Paired System-C and System-D panels display each operation, compact
state snapshot, CAS status, and components that differ from the oracle.  Selecting
the first highlighted row opens a deterministic explanation such as ``constraint
$x\neq 3$ is absent''; no LLM writes this diagnosis.

The recommended presentation mode is a controlled state-ownership comparison.  It
holds the CAS-constructed operation schedule and typed arguments fixed across the
two panels while the model still emits its complete claimed post-state.  System C
records that claim, whereas System D ignores the claim and re-derives state in its
ledger.  This removes operation-schedule differences from the visual comparison and
makes the manipulated variable explicit.  An autonomous diagnostic mode, in which
the model chooses every operation, remains available and labels schedule divergence
separately from stale-state drift.  In controlled mode, a premature final-answer
claim is recorded as state drift but does not terminate the fixed schedule.

The counterfactual sandbox operates on the trace just produced.  A user may edit an
operation or its arguments at any retained step, or edit the claimed state for
System C.  The same model can continue from the edited state, or the backend can
re-derive the downstream trace deterministically.  Every result is stored as a new
branch that can itself be edited again.  Failed operations are transactionally
rolled back to the last safe ledger state and retained as editable nodes; if a
protocol failure omitted a step, the aligned oracle operation provides a recovery
template.  A branch navigator and quick-restore controls support repeated sequential
interventions for either system.  System D continues to reject direct state edits
because its ledger owns state.  Sessions and individual branches
can be exported as JSON or Markdown with model configuration, operations, states,
metrics, intervention history, and failure events.  A labeled mock mode can plant a
chosen drift offline for tutorials; it is not used as experimental evidence.

\begin{figure}[t]
\centering
\IfFileExists{figures/statescope_ui.png}{%
  \includegraphics[width=\columnwidth]{figures/statescope_ui.png}%
}{%
  \fbox{\parbox[c][4.0cm][c]{.92\columnwidth}{\centering
  \textbf{REPLACE BEFORE SUBMISSION}\\[4pt]
  Add a high-resolution screenshot as\\
  \texttt{figures/statescope\_ui.png}.\\[6pt]
  Show both system panels, the highlighted first-drift row, and the drift-analysis
  drawer.}}%
}
\caption{The intended demo view: paired model-owned and ledger-owned traces,
with the first divergent state highlighted and explained.}
\label{fig:ui}
\end{figure}

\section{Framework Validation}

\subsection{Data and setup}

To validate the complete framework rather than position leaderboard numbers as its
main contribution, we instantiate DemoBench with 250 synthetic problems, balanced
across four families:
(A) binding chains and substitution; (B) equations with branch or domain
constraints; (C) recurrences and iterative state; and (D) derivative lemma DAGs.
Family counts are 63/63/62/62.  Every instance has fresh parameters, explicit
provenance, and a SymPy-constructed gold trace.  The benchmark therefore measures the
controlled phenomenon rather than broad contest-math ability, and avoids direct
test-set contamination from established math corpora.

We evaluate Qwen2.5-1.5B-Instruct \citep{qwen2024qwen25}, Qwen3-4B, Qwen3-8B, and
Qwen3-30B-A3B-Instruct-2507 \citep{yang2025qwen3}.  Each model is served locally
with vLLM in bfloat16 and temperature zero.  For every model--problem pair we run
Systems C and D once through the same text--JSON protocol.  This yields 2,000
completed runs (four models $\times$ 250 problems $\times$ two systems).  SymPy is
the sole judge.  Results are descriptive for this fixed, single-sample setting.

\begin{table*}[t]
\centering
\small
\begin{tabular}{lrrrrrr}
\toprule
& \multicolumn{2}{c}{Final accuracy (\%)}
& \multicolumn{2}{c}{$\mathrm{SF}_{g}$ (\%)}
& \multicolumn{2}{c}{Emitted-step coverage (\%)} \\
\cmidrule(lr){2-3}\cmidrule(lr){4-5}\cmidrule(lr){6-7}
Model & C & D & C & D & C & D \\
\midrule
Qwen2.5-1.5B          & 0.4 & 6.8  & 0.0 & 11.7 & 82.5 & 65.1 \\
Qwen3-4B              & 6.8 & 23.2 & 5.0 & 24.4 & 89.5 & 70.1 \\
Qwen3-8B              & 6.8 & 25.6 & 3.2 & 28.7 & 95.5 & 80.2 \\
Qwen3-30B-A3B         & 9.2 & 38.4 & 9.1 & 32.0 & 93.4 & 79.7 \\
\midrule
Pooled                 & 5.8 & 23.5 & 4.3 & 24.2 & 90.2 & 73.8 \\
\bottomrule
\end{tabular}
\caption{Paired evaluation on 250 problems per model.  C records model-claimed
state; D records the typed ledger.  $\mathrm{SF}_{g}$ penalizes missing steps as in
Eq.~\ref{eq:sfg}.  Coverage clarifies the trade-off introduced when D stops on an
invalid operation.}
\label{tab:main}
\end{table*}

\subsection{Results}

Table~\ref{tab:main} is a framework sanity check and an example of the controlled
comparisons \systemname{} enables.  It shows a consistent ledger benefit.  Pooled final-answer
accuracy rises by 17.7 points, and strict state fidelity rises by 19.9 points.  The
gain appears for every evaluated model rather than only the largest.  The ledger is
more conservative: emitted-step coverage falls from 90.2\% to 73.8\% because it
halts on an invalid or unverifiable operation.  This is desirable for auditability
but means that the accuracy gain should not be described as unconstrained solving
performance.

Final-answer checking also misses useful trace information.  In the model-owned
condition, 58 of 1,000 runs have a correct final answer, yet all 58 diverge from the
canonical state trace at an earlier aligned step.  Across both conditions, first
divergence most often involves bindings (39.7\% of component flags), followed by
constraints (20.6\%) and the current expression (18.4\%).  These are precisely the
components surfaced by the UI.  The result supports \systemname{} as an inspection
tool, while the high overall drift rates also motivate conservative wording:
trajectory mismatch is evidence to inspect, not by itself a proof that every
alternative derivation is invalid.

\section{Related Systems}

Outcome verifiers rank complete solutions \citep{cobbe2021verifiers}, while process
supervision assigns step-level feedback \citep{lightman2023verify}.  ProcessBench
asks models to locate the earliest incorrect step in fixed solutions
\citep{zheng2025processbench}.  \systemname{} instead provides an extensible
execution framework: it dynamically constructs controlled instances, tracks
structured state online, uses exact symbolic comparison where supported, and allows
the user to replay an edited operation.

ReasonGraph visualizes sequential and tree-structured inference paths
\citep{li2025reasongraph}.  Its focus is organizing reasoning text and method-level
paths; ours is semantic state comparison and CAS-backed diagnosis in mathematics.
Formal proof assistants expose proof states but require proofs in a formal language.
\systemname{} occupies a lighter-weight middle ground: a constrained operation
language, a computer-algebra oracle, and a browser interface for open language
models.

\section{Availability, Licensing, and Reproducibility}

The application is a Python backend with a dependency-free browser frontend.  Model
weights remain local and are served through vLLM; the supported evaluation path does
not call proprietary APIs.  The repository includes model configurations, dynamic
dataset generation, resumable batch execution, append-only traces, report generation,
and offline tests.  The software release is intended under \softwarelicense; this field
and the repository's \texttt{LICENSE} file must agree before submission.  Synthetic
DemoBench items are marked CC0-1.0, and upstream model licenses remain their owners'.
The public package/install page is \url{\demourl}.

\section{Limitations and Ethical Considerations}

Dynamic generation reduces instance-memorization risk but cannot establish that a
model has never seen the underlying templates.  DemoBench is synthetic and covers
algebraic state phenomena, not the breadth of
natural-language or competition mathematics.  Each model is sampled once, so we do
not estimate decoding variance.  Step-index alignment assumes the controlled
operation schedule and can flag a valid alternative derivation; CAS status and the
strict metric reduce but do not eliminate this construct-validity risk.  SymPy
parsing can also fail on informal notation.  Finally, the current evaluation tests
system behavior, not whether the interface helps people diagnose errors faster; a
counterbalanced raw-trace versus UI study is future work.

The system should not be used to certify high-stakes mathematical conclusions.
Its deterministic explanations describe differences from a reference trace, and
the UI preserves CAS and parser uncertainty rather than presenting every highlight
as a proof of error.  All headline data are freshly generated synthetic instances;
no personal or sensitive data are used.

\section{Conclusion}

\systemname{} is a framework for building controlled, inspectable reasoning
experiments rather than only another static math benchmark.  Dynamic generation
reduces contamination risk; explicit state makes context controllable; iterative
execution and counterfactual replay make reasoning genuinely interactive; and
core registries allow new families, model backends, and systems to reuse the same
runtime, while operations derive from one specification table.  Solution-state
drift is the flagship demonstration: contrasting
model-owned state with a CAS-verified ledger localizes divergence and makes the
consequence of an edit visible.  The demo is intended for researchers developing
reasoning agents, educators teaching multi-step derivations, and developers
evaluating tool-using models.

\begin{thebibliography}{8}

\bibitem[Cobbe et~al.(2021)Cobbe, Kosaraju, Bavarian, Chen, Jun, Kaiser,
Plappert, Tworek, Hilton, Nakano, Hesse, and Schulman]{cobbe2021verifiers}
Karl Cobbe, Vineet Kosaraju, Mohammad Bavarian, Mark Chen, Heewoo Jun, Lukasz
Kaiser, Matthias Plappert, Jerry Tworek, Jacob Hilton, Reiichiro Nakano,
Christopher Hesse, and John Schulman. 2021.
\newblock Training verifiers to solve math word problems.
\newblock \emph{arXiv preprint arXiv:2110.14168}.

\bibitem[Li et~al.(2025)Li, Shareghi, and Collier]{li2025reasongraph}
Zongqian Li, Ehsan Shareghi, and Nigel Collier. 2025.
\newblock ReasonGraph: Visualization of reasoning methods and extended inference
paths.
\newblock In \emph{Proceedings of ACL 2025: System Demonstrations}, pages 140--147.

\bibitem[Lightman et~al.(2023)Lightman, Kosaraju, Burda, Edwards, Baker, Lee,
Leike, Schulman, Sutskever, and Cobbe]{lightman2023verify}
Hunter Lightman, Vineet Kosaraju, Yura Burda, Harri Edwards, Bowen Baker, Teddy
Lee, Jan Leike, John Schulman, Ilya Sutskever, and Karl Cobbe. 2023.
\newblock Let's verify step by step.
\newblock \emph{arXiv preprint arXiv:2305.20050}.

\bibitem[Meurer et~al.(2017)Meurer, Smith, Paprocki, \v{C}ert\'{i}k,
Kirpichev, Rocklin, Kumar, Ivanov, Moore, Singh, et~al.]{meurer2017sympy}
Aaron Meurer, Christopher~P. Smith, Mateusz Paprocki, Ond\v{r}ej \v{C}ert\'{i}k,
Sergey~B. Kirpichev, Matthew Rocklin, AMiT Kumar, Sergiu Ivanov, Jason~K. Moore,
Sartaj Singh, et~al. 2017.
\newblock SymPy: Symbolic computing in Python.
\newblock \emph{PeerJ Computer Science}, 3:e103.

\bibitem[Qwen et~al.(2024)]{qwen2024qwen25}
Qwen, An Yang, Baosong Yang, Beichen Zhang, Binyuan Hui, Bo Zheng, Bowen Yu,
et~al. 2024.
\newblock Qwen2.5 technical report.
\newblock \emph{arXiv preprint arXiv:2412.15115}.

\bibitem[Yang et~al.(2025)Yang, Li, Yang, Zhang, Hui, Zheng, Yu, et~al.]{yang2025qwen3}
An Yang, Anfeng Li, Baosong Yang, Beichen Zhang, Binyuan Hui, Bo Zheng, Bowen Yu,
et~al. 2025.
\newblock Qwen3 technical report.
\newblock \emph{arXiv preprint arXiv:2505.09388}.

\bibitem[Zheng et~al.(2025)Zheng, Zhang, Zhang, Lin, Lu, Yu, Liu, Zhou, and
Lin]{zheng2025processbench}
Chujie Zheng, Zhenru Zhang, Beichen Zhang, Runji Lin, Keming Lu, Bowen Yu,
Dayiheng Liu, Jingren Zhou, and Junyang Lin. 2025.
\newblock ProcessBench: Identifying process errors in mathematical reasoning.
\newblock In \emph{Proceedings of ACL 2025}, pages 1009--1024.

\end{thebibliography}

\end{document}
